\section{Introdução}

O Problema do Caixeiro Viajante (TSP) pede o ciclo hamiltoniano de menor
custo num grafo completo com pesos. Ele é NP-difícil: algoritmos exatos
levam tempo exponencial no pior caso, e na prática são usados
aproximativos com garantia de qualidade quando a métrica satisfaz a
desigualdade triangular~\cite{applegate2006}. São eles o
Twice-around-the-tree (TATT), com fator 2, e o
Christofides~\cite{christofides1976}, com fator 1,5. Este trabalho
implementa três algoritmos e os avalia empiricamente:
\begin{itemize}
  \item um Branch-and-Bound (B\&B) exato com três limites inferiores (duas
    menores arestas, 1-árvore e Held-Karp~\cite{heldkarp1970,heldkarp1971});
  \item o TATT;
  \item o Christofides.
\end{itemize}
O foco não é só a escrita do código, mas também a \emph{auditoria} e a
\emph{modelagem empírica} do desempenho.

As contribuições são quatro.
\begin{enumerate}
  \item \emph{Profiling por subrotina e complexidade empírica}
    (Seção~\ref{sec:complexidade}). Mostramos onde a complexidade teórica
    aparece (o blossom, $N^{2{,}89}$) e onde ela é mascarada pelo
    interpretador (Prim $O(n^2)$ medido como $N^{1{,}2}$). Também mostramos
    que o garbage collector não distorce as curvas, mas cópias de memória
    sim.
  \item \emph{Impacto do warm-start} nos nós expandidos e podados do B\&B
    (Seção~\ref{sec:warmstart}). O ganho depende da ordem de visita dos
    filhos: quando a busca é cega, o fecho convexo reduz a árvore em
    41$\times$; quando os filhos são ordenados pelo limite, o UB inicial
    quase não importa.
  \item \emph{Análise de sensibilidade} (Seção~\ref{sec:sens}). O gap de
    cada aproximativo foi relacionado a features geométricas e da MST em
    174 instâncias com ótimo exato, 159 delas sintéticas e geradas para
    serem adversariais.
  \item \emph{Auditoria} (Seção~\ref{sec:auditoria}). Comparamos o código
    refatorado com uma versão ingênua e registramos uma revisão cruzada
    entre dois modelos de IA, que encontrou erros
    de \emph{medição}, e não de algoritmo.
\end{enumerate}
